%%
% BIThesis 研究生学位论文模板 The BIThesis Template for Graduate Thesis
% This file has no copyright assigned and is placed in the Public Domain.
%%

\chapter{系统总体框架与相关理论基础}

\label{chap:framework_theory}

\section{引言}

矿井巷道场景通常具有空间狭长、结构重复、纹理稀缺以及卫星导航不可用等特点，车辆位姿估计与环境建图难以依赖单一传感器独立完成。为了在此类场景中实现稳定的环境感知与定位，需要综合利用激光雷达提供的三维几何观测、惯性测量单元提供的高频运动信息以及由先验地图提取的结构约束信息，在统一框架下完成状态估计、地图构建与在线定位。

本文的整体研究工作围绕“离线建图形成先验、在线定位利用先验”这一主线展开。其中，第 3 章面向矿井巷道建图问题，重点研究如何利用几何结构约束提升先验地图的全局一致性；第 4 章则面向在线定位问题，重点研究如何结合先验体素约束与轴向分解机制，提高矿井场景下的定位稳定性。为便于后续章节展开，本章首先对系统总体框架进行说明，并对本文涉及的三维空间运动描述、点云配准与状态估计、图优化与因子图等相关理论基础进行统一介绍。

\section{系统总体框架}

本文研究的矿井巷道感知与定位系统总体上可分为离线建图与在线定位两个阶段。离线阶段以多源传感器采集的数据为基础，通过点云配准、位姿图优化以及几何结构约束建模，构建全局一致的先验地图，并进一步提取后续在线定位所需的先验体素表、巷道中轴线及相关结构信息。在线阶段则以当前时刻激光点云和 IMU 数据为输入，在先验地图支持下完成位姿预测、局部配准、全局配准以及差异化融合估计，从而获得当前车辆位姿。

从功能层次上看，该系统可进一步概括为数据获取层、先验构建层和在线定位层。数据获取层负责采集激光雷达点云、IMU 测量以及必要的辅助位姿信息；先验构建层对应本文第 3 章的研究内容，主要输出全局先验地图及结构几何先验；在线定位层对应本文第 4 章的研究内容，主要利用当前观测和先验信息进行递进式位姿估计与鲁棒融合。通过将离线建图与在线定位有机衔接，系统既能够利用全局先验增强定位稳定性，又能够保持对局部动态观测变化的适应能力。

[此处插入图 2.1 系统总体框架图]

建议图 2.1 采用“三层结构”的流程图形式绘制：
1. 左侧为输入层，包括激光雷达点云、IMU 数据以及可选的外部位姿信息；
2. 中间为处理层，分为离线建图支路和在线定位支路；
3. 右侧为输出层，包括先验地图、先验体素表、巷道中轴线以及在线位姿估计结果。

[此处插入表 2.1 系统主要模块及输入输出关系表]

表 2.1 可列出“模块名称--输入--输出--对应章节”四列，以便清晰说明全文结构与系统功能之间的对应关系。例如，离线建图模块可对应第 3 章，在线定位模块可对应第 4 章。

\section{三维空间运动描述}

为了统一描述矿井车辆、传感器与先验地图之间的几何关系，本文采用三维空间中的刚体运动模型表示系统状态，并在统一坐标系定义下描述不同观测之间的变换关系。该部分内容为后续点云配准、状态估计与图优化建模提供基本的数学表示。

\subsection{坐标系定义}

结合本文所采用的激光雷达与 IMU 组合系统，可定义以下几类常用坐标系：

1. 全局坐标系 $\{G\}$：用于表示先验地图及车辆全局位姿的参考坐标系。在离线建图阶段，优化后的全局地图通常建立在该坐标系下；在线定位阶段，车辆位姿也以该坐标系为基准进行表达。
2. 车体坐标系 $\{B\}$：固定于车辆本体，通常取车辆几何中心或安装基准点为原点，用于统一表达各传感器相对于车辆本体的外参关系。
3. 激光雷达坐标系 $\{L\}$：以激光雷达本体为参考建立的坐标系，原始点云数据通常在该坐标系下输出。
4. IMU 坐标系 $\{I\}$：以惯性测量单元安装位置为参考建立的坐标系，加速度计和陀螺仪测量均在该坐标系下给出。

在系统标定完成后，坐标系 $\{L\}$、$\{I\}$ 与 $\{B\}$ 之间的相对位姿关系可视为常量，而车辆在全局坐标系 $\{G\}$ 下的位姿则随时间变化。通过统一的坐标系定义，来自不同传感器的观测可被映射到同一参考系中，从而便于后续配准、融合与优化处理。

[此处插入图 2.2 系统坐标系定义示意图]

图 2.2 建议绘制车辆、激光雷达和 IMU 的相对安装位置，并明确标注 $\{G\}$、$\{B\}$、$\{L\}$ 和 $\{I\}$ 四个坐标系之间的关系。

[此处插入表 2.2 系统主要坐标系及其物理意义]

\subsection{刚体运动表示}

在三维空间中，车辆或传感器的位姿通常由旋转与平移共同决定。本文采用刚体变换群 $SE(3)$ 表示任意两个坐标系之间的位姿关系。设坐标系 $\{B\}$ 到坐标系 $\{A\}$ 的刚体变换记为 $\mathbf{T}_{A}^{B}$，则有：
\begin{equation}
\mathbf{T}_{A}^{B}
=
\begin{bmatrix}
\mathbf{R}_{A}^{B} & \mathbf{t}_{A}^{B}\\
\mathbf{0}^{\mathrm T} & 1
\end{bmatrix}
\in SE(3)
\end{equation}

其中，$\mathbf{R}_{A}^{B}\in SO(3)$ 表示旋转矩阵，$\mathbf{t}_{A}^{B}\in \mathbb{R}^{3}$ 表示平移向量。

对于旋转部分，本文同时采用旋转矩阵与四元数两种形式进行描述。旋转矩阵具有表达直观、便于组合变换等优点，而四元数则具有参数紧凑、便于插值以及避免欧拉角奇异性的优势。设四元数记为：
\begin{equation}
\mathbf{q}
=
\left[
q_w,q_x,q_y,q_z
\right]^{\mathrm T}
\end{equation}

则其对应的旋转矩阵可记为 $\mathbf{R}(\mathbf{q})$。在后续状态估计与配准建模中，本文通常以四元数表示姿态变量，并在需要构造刚体变换矩阵时将其转换为旋转矩阵形式。

若在连续两个时刻 $i$ 和 $j$ 分别得到系统位姿 $\mathbf{T}_i$ 和 $\mathbf{T}_j$，则两者之间的相对位姿可表示为：
\begin{equation}
\Delta \mathbf{T}_{i,j}
=
\mathbf{T}_{i}^{-1}\mathbf{T}_{j}
\end{equation}

该相对变换描述了系统从时刻 $i$ 到时刻 $j$ 的刚体运动关系，是后续帧间配准、里程计约束以及图优化建模中的基本量。

为便于描述位姿扰动和优化误差，常将刚体变换进一步映射到李代数空间 $\mathfrak{se}(3)$ 中。设位姿扰动向量记为 $\boldsymbol{\xi}\in \mathbb{R}^{6}$，则可通过指数映射与对数映射在 $SE(3)$ 与 $\mathfrak{se}(3)$ 之间进行转换，即：
\begin{equation}
\mathbf{T}
=
\mathrm{Exp}\left(\boldsymbol{\xi}^{\wedge}\right),
\qquad
\boldsymbol{\xi}
=
\mathrm{Log}\left(\mathbf{T}\right)^{\vee}
\end{equation}

其中，$(\cdot)^{\wedge}$ 与 $(\cdot)^{\vee}$ 分别表示向量与李代数矩阵之间的对应关系。该表示方式在非线性优化和误差状态建模中具有重要作用。

\subsection{坐标变换关系}

在多传感器系统中，来自不同坐标系下的点、向量或位姿需要通过坐标变换映射到统一参考系中。若一点 $\mathbf{p}^{B}$ 在坐标系 $\{B\}$ 下的齐次坐标表示为：
\begin{equation}
\bar{\mathbf{p}}^{B}
=
\left[
\left(\mathbf{p}^{B}\right)^{\mathrm T},1
\right]^{\mathrm T}
\end{equation}

则其在坐标系 $\{A\}$ 下的坐标可表示为：
\begin{equation}
\bar{\mathbf{p}}^{A}
=
\mathbf{T}_{A}^{B}\bar{\mathbf{p}}^{B}
\end{equation}

若考虑激光雷达点云从传感器坐标系映射到全局坐标系的过程，则对于激光点 $\mathbf{p}^{L}$，其在全局坐标系 $\{G\}$ 下的坐标满足：
\begin{equation}
\bar{\mathbf{p}}^{G}
=
\mathbf{T}_{G}^{B}\mathbf{T}_{B}^{L}\bar{\mathbf{p}}^{L}
\end{equation}

其中，$\mathbf{T}_{B}^{L}$ 为激光雷达相对于车体的外参变换，$\mathbf{T}_{G}^{B}$ 为车体在全局坐标系下的位姿。若系统以 IMU 作为主体状态参考，则也可写为：
\begin{equation}
\bar{\mathbf{p}}^{G}
=
\mathbf{T}_{G}^{I}\mathbf{T}_{I}^{L}\bar{\mathbf{p}}^{L}
\end{equation}

该式表明，点云的全局位置由车辆全局位姿与传感器外参共同决定。

此外，坐标变换还满足链式组合关系与逆变换关系，即：
\begin{equation}
\mathbf{T}_{A}^{C}
=
\mathbf{T}_{A}^{B}\mathbf{T}_{B}^{C}
\end{equation}

\begin{equation}
\mathbf{T}_{B}^{A}
=
\left(\mathbf{T}_{A}^{B}\right)^{-1}
=
\begin{bmatrix}
\left(\mathbf{R}_{A}^{B}\right)^{\mathrm T} &
-\left(\mathbf{R}_{A}^{B}\right)^{\mathrm T}\mathbf{t}_{A}^{B}\\
\mathbf{0}^{\mathrm T} & 1
\end{bmatrix}
\end{equation}

这些基本关系为后续点云配准中的位姿初始化、传感器测量对齐以及状态估计中的观测方程构建提供了统一的数学基础。

\section{点云配准与状态估计基础}

点云配准的目标是在两个或多个点云观测之间求解最优刚体变换，使得不同观测下的几何结构尽可能对齐。按照观测对象不同，点云配准可分为帧间配准和帧到地图配准两类：前者主要用于估计局部连续运动，后者主要用于建立当前观测与全局先验之间的几何联系。无论采用何种具体算法，其本质均可抽象为在给定对应关系或局部结构约束条件下，对刚体变换参数进行优化求解。

设源点云为 $\mathcal{P}=\{\mathbf{p}_i\}$，目标点云为 $\mathcal{Q}=\{\mathbf{q}_i\}$，待估计刚体变换为 $\mathbf{T}\in SE(3)$，则一般形式的点云配准目标函数可表示为：
\begin{equation}
\hat{\mathbf{T}}
=
\arg \min_{\mathbf{T}\in SE(3)}
\sum_{i\in \mathcal{C}}
\mathbf{e}_{i}^{\mathrm T}
\boldsymbol{\Omega}_{i}
\mathbf{e}_{i}
\end{equation}

其中，$\mathcal{C}$ 表示对应关系集合，$\boldsymbol{\Omega}_{i}$ 为权重或信息矩阵，$\mathbf{e}_{i}$ 表示在当前变换下构造的匹配残差。根据残差定义方式不同，可进一步得到点到点、点到平面、分布匹配等不同形式的配准模型。对于矿井巷道场景而言，由于存在局部平面结构、结构重复以及点云稀疏程度变化等问题，配准算法的选择与初值质量对最终估计结果具有重要影响。

在实际系统中，激光点云通常以较低频率提供高精度几何观测，而 IMU 则以较高频率提供短时运动信息。因此，状态估计问题通常采用“高频预测、低频校正”的方式进行求解。设系统在时刻 $k$ 的状态记为 $\mathbf{x}_{k}$，控制输入记为 $\mathbf{u}_{k}$，观测记为 $\mathbf{z}_{k}$，则贝叶斯滤波框架下的状态估计过程可概括为预测与更新两个步骤：
\begin{equation}
p\left(\mathbf{x}_{k}\mid \mathbf{z}_{1:k-1},\mathbf{u}_{1:k}\right)
=
\int
p\left(\mathbf{x}_{k}\mid \mathbf{x}_{k-1},\mathbf{u}_{k}\right)
p\left(\mathbf{x}_{k-1}\mid \mathbf{z}_{1:k-1},\mathbf{u}_{1:k-1}\right)
\mathrm{d}\mathbf{x}_{k-1}
\end{equation}
\begin{equation}
p\left(\mathbf{x}_{k}\mid \mathbf{z}_{1:k},\mathbf{u}_{1:k}\right)
\propto
p\left(\mathbf{z}_{k}\mid \mathbf{x}_{k}\right)
p\left(\mathbf{x}_{k}\mid \mathbf{z}_{1:k-1},\mathbf{u}_{1:k}\right)
\end{equation}

其中，第一式表示根据上一时刻状态和当前控制输入进行先验传播，第二式表示根据当前观测对先验状态进行修正。在本文研究中，IMU 测量主要承担运动预测作用，而点云配准结果主要承担观测校正作用。通过将二者有机结合，可同时兼顾状态估计的连续性与几何约束能力。

需要指出的是，点云配准与状态估计在系统中并非相互独立。高质量的状态预测能够为配准提供更好的初始位姿，从而改善配准收敛性；而稳定的配准结果又能够反过来修正状态预测误差。因此，在矿井巷道场景下，如何构建合理的配准观测，并与运动预测结果进行有效融合，是提升系统定位稳定性的关键基础。

\section{图优化与因子图基础}

对于需要利用多时刻、多类型观测进行联合估计的问题，图优化提供了一种统一而有效的建模框架。其基本思想是：将待估计的状态变量表示为图中的节点，将观测约束表示为连接状态变量的边或因子，从而将复杂的联合估计问题转化为一个带有多类约束的非线性最小二乘优化问题。该思想在位姿图优化、后端建图以及多源几何约束联合求解中均具有广泛应用。

在因子图表示中，设状态变量集合为：
\begin{equation}
\mathcal{X}
=
\left\{
\mathbf{x}_{1},\mathbf{x}_{2},\ldots,\mathbf{x}_{n}
\right\}
\end{equation}

观测集合为：
\begin{equation}
\mathcal{Z}
=
\left\{
\mathbf{z}_{1},\mathbf{z}_{2},\ldots,\mathbf{z}_{m}
\right\}
\end{equation}

则在条件独立假设下，后验概率分布可写为多个因子的乘积形式，即：
\begin{equation}
p\left(\mathcal{X}\mid \mathcal{Z}\right)
\propto
\prod_{j=1}^{m}
\phi_{j}\left(\mathcal{X}_{j}\right)
\end{equation}

其中，$\phi_{j}(\mathcal{X}_{j})$ 表示与观测 $\mathbf{z}_{j}$ 对应的因子，$\mathcal{X}_{j}$ 表示该因子所关联的状态变量子集。通过这种表示方式，不同来源的约束信息可以被统一组织到同一图结构中。

当进一步假设各类观测误差服从高斯分布时，最大后验估计问题可等价转化为残差加权平方和最小化问题，即：
\begin{equation}
\mathcal{X}^{\ast}
=
\arg \min_{\mathcal{X}}
\sum_{j=1}^{m}
\mathbf{e}_{j}^{\mathrm T}
\boldsymbol{\Omega}_{j}
\mathbf{e}_{j}
\end{equation}

其中，$\mathbf{e}_{j}$ 表示第 $j$ 个因子的残差项，$\boldsymbol{\Omega}_{j}$ 表示对应的信息矩阵。该目标函数形式与配准问题中的残差优化具有一致性，因此图优化可以看作是对多类观测残差进行统一求解的推广形式。

在实际求解中，由于残差函数通常关于状态变量是非线性的，需要在当前估计点附近对其进行一阶线性化。设状态增量记为 $\delta \mathcal{X}$，则残差可近似表示为：
\begin{equation}
\mathbf{e}_{j}\left(\mathcal{X}\oplus \delta \mathcal{X}\right)
\approx
\mathbf{e}_{j}\left(\mathcal{X}\right)
+\mathbf{J}_{j}\delta \mathcal{X}
\end{equation}

其中，$\mathbf{J}_{j}$ 为残差对状态变量的雅可比矩阵，$\oplus$ 表示状态更新算子。将全部因子线性化后，可得到标准的正规方程：
\begin{equation}
\mathbf{H}\delta \mathcal{X}
=
\mathbf{b}
\end{equation}

其中，$\mathbf{H}$ 为近似海森矩阵，$\mathbf{b}$ 为梯度向量。随后可通过高斯牛顿法、列文伯格--马夸尔特法等非线性优化方法迭代求解最优状态。

[此处插入图 2.3 因子图基本结构示意图]

图 2.3 建议采用“变量节点 + 因子节点”的典型结构示意图，展示相邻位姿变量、里程计因子、先验因子以及几何约束因子之间的连接关系。该图可为第 3 章中位姿图优化建模提供概念铺垫。

从本文的研究任务出发，图优化与因子图理论主要服务于离线建图阶段的后端统一求解。第 3 章中，多种几何约束将被表示为因子图中的不同因子，并在统一框架下对关键帧位姿进行联合优化。通过在第二章先给出通用的理论形式，可以避免后续章节在进入具体建模时重复展开基础概念。

\section{建图结果评价方法}

SLAM 系统的性能评价通常包括轨迹精度评价和地图质量评价两个方面。对于具有真值轨迹的数据集，常采用绝对位姿误差（Absolute Pose Error, APE）等指标衡量估计轨迹与参考轨迹之间的偏差，从而定量反映系统在全局尺度上的定位与建图精度。该类指标定义明确、可解释性强，适用于公开数据集或仿真数据集中能够获得高精度真值轨迹的情形。

然而，在真实矿井等工程场景中，长距离、高精度真值轨迹往往难以获取，此时仅依赖轨迹误差难以全面评价建图结果。针对这一问题，近年来出现了面向点云地图质量评估的相关工作，其中 MapEval 从全局几何一致性和局部结构一致性两个方面对 SLAM 地图进行统一评价。当存在高精度参考地图时，可将待评估地图与参考地图进行配准，并进一步从地图层面对不同算法的建图结果进行定量比较。因此，本文在具有真值轨迹的公开数据集和仿真数据集上采用 APE 评价轨迹精度，同时结合地图对比结果进行分析；在缺少真值轨迹但具有参考地图的真实地下矿井数据集上，则采用基于参考地图的 MapEval 评价方法对建图质量进行补充评估。

\section{本章小结}

本章对本文研究所依赖的系统总体框架与相关理论基础进行了统一说明。首先，从离线建图与在线定位两个阶段出发，介绍了矿井巷道感知与定位系统的整体组成及各模块之间的关系。随后，围绕三维空间运动描述，对常用坐标系、刚体运动表示以及坐标变换关系进行了定义，为后续位姿建模提供统一的数学基础。接着，对点云配准与状态估计的基本思想进行了概述，说明了几何观测与运动预测在系统中的互补作用。最后，对图优化与因子图的基本建模方式及求解过程进行了介绍，为后续章节中的位姿图优化与多因子联合建模奠定了理论基础。

在此基础上，下一章将进一步面向矿井巷道先验地图构建问题，研究如何利用几何结构约束提升建图结果的全局一致性与结构准确性。
